\MakeAbstract{
    本课程设计围绕《深度学习模块：手写数字识别与垃圾分类》的要求展开。实验首先使用 MindSpore 在 MNIST 手写数字数据集上从零训练 LeNet5，并保存预训练权重；随后将该权重加载到 26 类垃圾分类数据集上，完成从零训练、冻结微调、全量微调、LoRA 微调和模型改进实验。垃圾分类图片在输入 LeNet5 前统一进行 RGB 转灰度、resize 到 $32\times32$、归一化和 HWC 到 CHW 的格式转换。在完成基础迁移实验后，本文进一步使用 MobileNetV2 进行进阶对照，分析更强骨干网络、RGB 输入、高分辨率和 ImageNet 预训练权重对垃圾分类性能的影响。

    实验代码包含数据加载、LeNet5 模型定义、预训练权重加载、冻结策略控制、LoRA 低秩适配、训练评估、CSV 日志记录、JSON 结果汇总和 PNG 曲线绘制等模块。冻结微调中冻结 \texttt{conv1}、\texttt{conv2}、\texttt{fc1}、\texttt{fc2}，只训练最后分类层 \texttt{fc3}；全量微调中所有层参与训练；LoRA 微调中在 \texttt{fc1} 和 \texttt{fc2} 上加入低秩分支，冻结原始主干权重，只训练 LoRA 参数和 \texttt{fc3}。

    实验结果表明：LeNet5 在 MNIST 上达到 98.60\% 的测试准确率；迁移到垃圾分类任务后，从零训练准确率为 3.85\%，冻结微调为 15.38\%，全量微调为 64.23\%，LoRA 微调为 25.38\%，强数据增强改进实验为 38.85\%。结果说明 MNIST 预训练权重可以作为有效初始化，但由于 MNIST 与垃圾分类图像分布差异较大，仅训练最后分类层的冻结微调适应能力不足；全量微调能够重新调整卷积层和全连接层，在 LeNet5 设置下取得最好效果。为进一步考察真实垃圾分类任务的性能上限，本文在正文中加入 MobileNetV2 进阶对照实验：MindCV 官方预训练 MobileNetV2-1.0 结合 warmup + cosine 全量微调达到 91.15\%，说明 RGB 输入、高分辨率和高质量 ImageNet 预训练对真实垃圾分类任务具有显著价值。
}{LeNet5；MNIST；垃圾分类；MindSpore；微调；LoRA}
